\section{Experimental Setup}

\subsection{Research Questions}

We formalize three research questions that structure our experimental design:

\textbf{RQ1:} Does the curated-corpus model (OLMo-7B / Dolma) achieve a higher MMLU/HellaSwag generalization balance ratio than the minimally-curated model (Pythia-6.9B / The Pile) at matched training scale ($\sim$300B tokens)?

\textbf{RQ2:} Do the two models differ in ARC difficulty delta ($\Delta_{\text{ARC}}$), indicating differences in reasoning robustness attributable to corpus quality?

\textbf{RQ3:} Is the HellaSwag denominator informative at this training scale, or does saturation collapse the ratio metric into a proxy for MMLU-only differences?

\subsection{Models and Checkpoints}

\begin{table}[h]
\centering
\caption{Model configurations and checkpoint details.}
\begin{tabular}{llllll}
\toprule
Model & Architecture & Parameters & HuggingFace ID & Revision & Tokens \\
\midrule
Pythia-6.9B & GPT-NeoX & 6.9B & \texttt{EleutherAI/pythia-6.9b} & \texttt{step143000} & $\approx$299.9B \\
OLMo-7B & OLMo & 7B & \texttt{allenai/OLMo-7B} & \texttt{step68000-tokens301B} & $\approx$301B \\
\bottomrule
\end{tabular}
\end{table}

Pythia-6.9B uses GPT-NeoX architecture with rotary positional embeddings and parallel attention/MLP layers. OLMo-7B uses the OLMo architecture with SwiGLU activations and a different tokenizer vocabulary (50,280 vs.\ 50,254 tokens). These architectural differences are a known confound discussed in Section~6.

\subsection{Evaluation Benchmarks}

\textbf{MMLU}~\cite{Hendrycks2021MMLU} (Massive Multitask Language Understanding) covers 57 academic subjects spanning STEM, humanities, social sciences, and professional domains. We report macro-averaged accuracy across all subjects using 5-shot prompting. MMLU serves as the numerator of our primary metric, capturing factual knowledge acquisition from pre-training data.

\textbf{HellaSwag}~\cite{Zellers2019HellaSwag} is a commonsense NLI benchmark requiring models to select the most plausible continuation of a given context from four adversarially-filtered options. We report normalized accuracy using 10-shot prompting. HellaSwag serves as the denominator of our primary metric, capturing commonsense generalization.

\textbf{ARC-Easy and ARC-Challenge}~\cite{Clark2018ARC} (AI2 Reasoning Challenge) provide multiple-choice science questions at two difficulty levels. ARC-Easy contains questions answerable by retrieval; ARC-Challenge contains questions requiring multi-step reasoning. We report accuracy using 25-shot prompting for both splits. The difficulty delta ($\Delta_{\text{ARC}}$) serves as our secondary metric.

\subsection{Evaluation Protocol}

All evaluations use \texttt{lm-evaluation-harness} v0.4.12~\cite{gao2024harness}. The following command reproduces the primary evaluation:

\begin{verbatim}
lm_eval --model hf \
  --model_args pretrained=EleutherAI/pythia-6.9b,revision=step143000 \
  --tasks mmlu,hellaswag,arc_easy,arc_challenge \
  --num_fewshot 5 \
  --limit 500 \
  --output_path results/pythia_6.9b_step143000 \
  --device cuda
\end{verbatim}

Replace \texttt{pretrained} and \texttt{revision} for the OLMo checkpoint. The \texttt{--limit 500} flag enables fast evaluation; remove for full benchmark evaluation. Hardware: NVIDIA H100 NVL. Estimated runtime per model: approximately 45 minutes with \texttt{--limit 500}.

\subsection{Statistical Analysis Protocol}

Bootstrap resampling uses $B = 1000$ iterations with replacement sampling over MMLU subjects. For each iteration, per-model ratios are computed from resampled subject accuracies, and the difference (OLMo $-$ Pythia) is recorded. The 95\% CI is the $[2.5, 97.5]$ percentile interval of the bootstrap distribution. Cohen's $d$ uses pooled standard deviation of per-subject MMLU accuracies as the scaling denominator. The hypothesis threshold is set at $+0.02$ (OLMo ratio exceeds Pythia by at least 0.02); refutation requires the CI to lie entirely below zero.
